% RSS 템플릿(IEEEtran) 한글 초안. 수업 과제 제출 후 영어로 옮겨 RSS 제출.
% Overleaf: 이 폴더(main.tex, IEEEtran.cls, references.bib)를 올리고 Menu > Compiler = XeLaTeX (kotex 필요).
% \todo{...} 는 아직 채우지 않은 부분. 최종본 전에 모두 지운다.
% 수치는 저장소의 실험 기록에서 그대로 옮긴 것. 출처는 각 줄 옆 주석(%)에 적었다.
% 범위(2026-11-27): 모든 실험은 작업자 오류가 있는 조건이다. 오류 없는 시험 조건은 넣지 않는다.
% 2026-11-28: 학습 18일 / 시험 처음 보는 6일로 설계 변경. 학습 날짜 과적합을 분석 절로 추가.
\documentclass[conference]{IEEEtran}
\usepackage{kotex}
\setmainfont{TeX Gyre Termes} % XeLaTeX용 Times. times 패키지는 XeLaTeX에서 글꼴을 못 찾는다.

% numbers option provides compact numerical references in the text.
\usepackage[numbers]{natbib}
\usepackage{multicol}
\usepackage{booktabs}
\usepackage{amsmath,amssymb}
\usepackage{xcolor}
\usepackage[bookmarks=true]{hyperref}
\newcommand{\todo}[1]{\textcolor{red}{[TODO: #1]}}

\begin{document}

\title{재고 기록이 어긋나는 물류창고에서의\\경량 시뮬레이터 기반 다중 로봇 작업 배정 학습}

% You will get a Paper-ID when submitting a pdf file to the conference system
\author{Author Names Omitted for Anonymous Review. Paper-ID [add your ID here]}

\maketitle

\begin{abstract}
물류창고의 창고 실행 시스템(WES)은 입고 적치, 출고 피킹, 재고 확인 작업을 여러 이동 로봇에게 실시간으로 배정한다.
현장 작업자가 기록 없이 재고를 옮기거나 빼 가면 시스템 기록과 실제 재고가 어긋나고, 로봇은 빈 칸에 가거나 확인 작업을 추가로 해야 한다.
본 연구는 이러한 기록 불일치가 있는 창고에서, 작업-로봇 짝에 선형 점수를 매기는 배정 정책의 가중치를 학습한다.
학습은 실제 WMS·WES 코드를 그대로 쓰고 로봇 이동만 단순화한 경량 시뮬레이터에서 교차 엔트로피 방법(CEM)으로 하고, 검증은 물리 시뮬레이터 Webots에서 한다.
실제 신발 제조사 창고의 주문 데이터에서 18일로 학습하고, 학습에 쓰지 않은 6일로 시험한다.
\todo{새 날짜 시험 결과(경량, Webots)로 주요 수치 작성}
또한 학습한 날짜를 시드만 바꿔 시험하면 개선이 크게 부풀려짐을 보인다. 학습 날짜에서 7.4\%였던 개선이 처음 보는 날짜에서는 1\% 안팎으로 줄었다.
\end{abstract}

\IEEEpeerreviewmaketitle

% ---------------------------------------------------------------
\section{서론}

자동화 물류창고에서는 여러 대의 이동 로봇이 입고된 상자를 랙에 적치하고(putaway), 주문에 맞춰 상자를 출고 레인으로 옮기며(pick),
기록과 실제가 다른지 의심되는 칸을 확인한다(verify)~\cite{wurman2008}.
어떤 작업을 어느 로봇에게 언제 맡길지 정하는 작업 배정은 창고의 처리 속도와 출고 마감 준수에 직접 영향을 준다~\cite{gerkey2004}.

실제 창고에서는 로봇과 사람이 함께 일한다.
작업자가 상자를 기록 없이 다른 칸으로 옮기거나 빼 가면 재고 기록이 틀어지고~\cite{dehoratius2008},
로봇은 상자가 없는 칸에 도착해 피킹에 실패하거나, 이미 차 있는 칸에 적치하려다 다른 칸을 찾아야 한다.
그 결과 확인 작업과 재작업이 생기고, 작업 대기열에는 피킹, 적치, 확인 작업이 섞인다.

배정 규칙은 보통 사람이 정한 우선순위(가까운 로봇, 피킹 우선 등)로 만든다.
이런 규칙은 기록이 정확하다는 가정에서 조정되었기 때문에, 확인 작업과 재작업이 섞이는 상황에서 어떤 작업을 먼저 할지를 잘 고르지 못한다.
본 연구는 기록 불일치가 있는 조건에서 배정 점수의 가중치를 데이터 기반 시나리오로 직접 학습하여 이 문제를 다룬다.

\textbf{기여.}
\begin{enumerate}
  \item 실제 창고 주문 데이터와 작업자 오류 모델로 만든 시나리오에서, 실제 WMS·WES 코드를 그대로 쓰는 경량 시뮬레이터로 배정 가중치를 학습하고 물리 시뮬레이터(Webots)에서 검증하는 절차를 제시한다.
  \item 학습에 쓰지 않은 날짜의 주문으로 학습한 배정을 팀 기본 WES, 거리 최적 매칭(헝가리안)과 비교한다. \todo{결과 요약 한 줄}
  \item 같은 날짜를 시드만 바꿔 시험하면 학습한 배정의 개선이 크게 부풀려짐을 보인다. 시드는 입고와 작업자 오류만 바꾸고 출고 주문은 그대로 두기 때문이다.
\end{enumerate}

% ---------------------------------------------------------------
\section{관련 연구}

\textbf{다중 로봇 작업 배정.}
\citet{gerkey2004}은 다중 로봇 작업 배정 문제를 분류했다.
\citet{choudhury2020}은 완료 시간의 불확실성과 시간 제약이 있는 동적 배정을 다뤘고,
\citet{dhanaraj2024}은 실패와 복구 작업을 고려한 hindsight optimization을 제안했다.
\todo{우리 방법과의 차이: 이들은 작업 실패 확률을 알고 계획하지만, 우리는 기록 불일치가 언제 어디서 생기는지 모르는 상태에서 WMS가 보여 주는 정보만으로 배정한다}

\textbf{물류창고 로봇.}
\citet{wurman2008}은 수백 대 규모 창고 로봇 시스템을 소개했고,
\citet{leet2023}은 작업 구성, 배정, 스케줄, 경로 계획을 결합해 대규모 작업량을 다뤘다.
본 연구는 그보다 작은 규모에서, 기록 불일치가 있는 조건의 배정 정책 학습에 집중한다.

\textbf{재고 기록 불일치.}
실제 재고와 시스템 기록이 다른 문제는 소매·물류에서 널리 보고되었다~\cite{dehoratius2008}.
검사와 보충을 결합한 정책~\cite{kok2007,akkerman2025}이 연구되었으나, 로봇 이동과 자원 경합을 함께 다루지는 않는다.

\textbf{정책 탐색.}
교차 엔트로피 방법(CEM)~\cite{rubinstein1999}은 경사 없이 정책 파라미터를 찾는 방법이다.
에피소드 끝에 한 번 나오는 비용과 적은 파라미터 수에 적합하다.

% ---------------------------------------------------------------
\section{문제 설정}

\subsection{창고와 작업}
창고는 $24\times14$\,m 크기이며 랙 160칸, 입고 대기칸(buffer), 지역별 출고 레인 4개(각 8칸)로 구성된다. 로봇은 8대이며, 모든 실험은 이 창고에서 한다. % scenario.json world, outbound.lanes
작업 종류는 다음과 같다.
\begin{itemize}
  \item \textbf{PUTAWAY}: 입고 대기칸의 상자를 정해진 랙 칸에 적치한다.
  \item \textbf{PICK}: 주문 품목이 있는 랙 칸의 상자를 해당 지역 출고 레인으로 옮긴다.
  \item \textbf{RELOCATE}: 상자를 다른 칸으로 옮긴다. 배정에서는 적치와 같은 종류로 다룬다.
  \item \textbf{VERIFY}: 칸의 실제 내용을 확인해 기록을 바로잡는다. 놀고 있는 로봇이 수행한다.
\end{itemize}
출고 트럭은 지역별로 900\,초마다 마감되어 출발한다. % cutoff_period_s
에피소드 길이는 3600\,초이다.

\subsection{시스템 구조}
창고 관리 시스템(WMS)은 재고 기록과 작업 상태를 관리하고 작업을 만든다.
WES는 WMS가 보여 주는 정보만으로 매 단계 작업과 로봇을 짝지어 실행 지시를 보낸다.
WES는 실제 재고나 미래 주문 같은 숨겨진 정보를 쓰지 않는다.
\todo{그림 1: WMS--WES--로봇 구조도}

\subsection{사람 작업자에 의한 기록 불일치}
현장 작업자가 기록 없이 재고를 옮기는 상황을 세 가지로 모델링한다.
\begin{itemize}
  \item S1: 상자를 다른 칸으로 옮긴다. 시간당 8회, 한 번에 최대 3칸이다.
  \item S2: 상자를 빼 간다. 시간당 3회이다.
  \item S3: 적치하려는 칸이 이미 차 있다. 확률 0.1이다.
\end{itemize}
불일치는 로봇이 칸에 도착해 라벨을 읽을 때 드러난다.
피킹 실패(SHORT\_PICK)나 품목 불일치(ITEM\_MISMATCH)가 나면 이웃 칸을, 적치 칸이 차 있으면(DESTINATION\_OCCUPIED) 그 칸을 확인 작업으로 만든다(열린 확인 작업 최대 4개). % sched_verify VERIFY_POLICY=event, VERIFY_MAX_OPEN 4
로봇 자체의 오류(적치 실수, 스캐너 실패)는 0으로 둔다.
모든 학습과 시험은 위 작업자 오류가 있는 조건에서 한다.
위 비율은 실측 자료가 없어 정한 합성 가정이며, 한계에서 다시 언급한다.

\subsection{시나리오}
주문은 공개된 신발 제조사 창고 피킹 데이터셋~\cite{assis2023}의 \texttt{Customer\_Order.csv}에서 하루를 골라 시간을 압축해 만든다.
\begin{itemize}
  \item 시간 압축 비율은 약 6배이고, 웨이브는 90\,초 간격으로 풀린다. % time_scale 6.08, wave_interval_s 90
  \item 입고는 그날 수요를 덮도록 합성한 트럭으로 만들고, 트럭 한 대에 상자 12개가 실린다.
  \item 날짜가 출고 주문(품목, 수량, 풀리는 시각)을 정한다. 날짜마다 한 시간 동안 다루는 상자 수는 약 120--2200개로 다르다.
  \item 시드는 입고 트럭의 도착 시각과 상자 구성, 작업자 오류가 언제 어디서 생기는지를 정한다. 출고 주문과 처음 재고 배치는 시드가 바뀌어도 같다.
\end{itemize}

\subsection{평가 지표}
\begin{itemize}
  \item \textbf{주문 리드타임}(주 지표): 주문이 풀린 시각부터 그 주문의 마지막 상자가 출하된 시각까지의 평균이다. 끝내지 못한 주문은 에피소드 종료 시각을 완료 시각으로 쳐서 포함한다. 끝낸 주문만 평균하면 어려운 주문을 미루는 정책이 유리해 보이기 때문이다.
  \item \textbf{상자 단위 리드타임}: 같은 방식으로 상자마다 잰 평균이다.
  \item \textbf{p90}: 주문 리드타임의 90번째 백분위수이다.
  \item \textbf{마감 놓침}: 마감까지 출하하지 못한 주문 수이다.
\end{itemize}

% ---------------------------------------------------------------
\section{방법}

\subsection{작업-로봇 짝 점수}
매 단계 시작 가능한 대기 작업 $\tau$ 각각을, 놀고 있는 로봇과 일하고 있는 로봇 $r$ 각각과 짝지어 점수를 매긴다.
\begin{equation}
  S(\tau, r) = \mathbf{w}^\top \boldsymbol{\phi}(\tau, r)
\end{equation}
점수가 높은 짝부터 탐욕적으로 고르며, 한 작업과 한 로봇은 한 번만 쓰인다.
놀고 있는 로봇이 이기면 바로 실행 지시를 보내고, 일하고 있는 로봇이 이기면 아무 지시도 보내지 않고 다음 단계에서 다시 정한다.
특징 $\boldsymbol{\phi}$는 표~\ref{tab:features}와 같다.
일하고 있는 로봇의 거리는 \emph{현재 작업의 남은 이동 거리 + 현재 작업이 끝나는 곳에서 새 작업까지의 거리}로 계산한다.
모든 특징은 WMS가 보여 주는 정보로만 계산한다.

\begin{table}[t]
\centering
\caption{배정 점수의 특징. 위 세 개만 로봇에 따라 달라진다.} % lib/sched_learned.py FEATURES
\label{tab:features}
\small
\begin{tabular}{@{}lp{0.62\columnwidth}@{}}
\toprule
특징 & 정의 \\
\midrule
dist & 로봇에서 작업 출발 칸까지 거리 (m) / 10 \\
busy & 일하고 있는 로봇이면 1 \\
eta & 로봇이 작업에 도착할 예상 시간 (분) \\
\midrule
pick, putaway, verify & 작업 종류 (0/1) \\
age & 대기 시간 / 600\,s (최대 2) \\
slack & 마감까지 남은 시간 / 900\,s, [0,1] (PICK) \\
truck\_here & 해당 레인에 트럭이 있으면 1 (PICK) \\
lane\_fill & 출고 레인 점유율 (PICK) \\
buffer\_fill & 입고 대기칸 점유율 (PUTAWAY) \\
\bottomrule
\end{tabular}
\end{table}

\subsection{경량 시뮬레이터}
Webots 실행은 1시간 시나리오에 약 8.5분이 걸려 학습에 쓰기 어렵다.
그래서 같은 시나리오 파일과 같은 WMS·WES 코드를 쓰고, 로봇 이동만 단순화한 경량 시뮬레이터를 만들었다.
작업자 오류도 같은 코드로 들어가므로, 학습 중에 기록 불일치와 확인 작업이 그대로 나타난다.
\begin{itemize}
  \item 한 구간의 이동 시간은 \emph{경로 길이 / 속도 + 2\,초 + 1\,초 $\times$ (우선권이 앞선 움직이는 로봇 수)}이다.
  \item 시간 간격은 0.5\,초이고, 충돌과 회피는 모델링하지 않는다.
  \item 1시간 시나리오를 약 1초에 실행한다.
\end{itemize}
Webots 실행 51개와 비교했을 때 가동률 평균 오차는 2.2\%, 주문 리드타임 평균 오차는 5.4\%였다. % docs/lite_sim.md, 2026-11-16
방법 간 순위는 리드타임 기준 92\%, 출하 기준 86\%가 일치했다.
\todo{이 보정은 Webots에 비켜서기 동작을 넣기 전 결과. 재보정하거나 한계에 명시}

\subsection{학습}
가중치 $\mathbf{w}$는 작업자 오류가 있는 경량 시뮬레이터에서 CEM으로 학습한다.
\begin{itemize}
  \item 후보 24개, 상위 6개, 15회 반복, 초기 표준편차 1.0(최소 0.05)이다.
  \item 시작값은 \emph{가까운 로봇, 피킹 우선} 규칙이다(dist $-1$, busy $-1$, pick $+2$).
  \item 매 반복 모든 후보를 학습 시나리오 36개 전부에서 평가해 같은 조건으로 비교한다.
  \item 학습 난수 시드를 3개(r0--r2) 바꿔 결과의 흔들림을 잰다.
\end{itemize}
에피소드 비용은 다음과 같다.
\begin{equation}
  C = 0.5\,L_{\text{order}} + 0.5\,L_{\text{unit}} + \lambda\,N_{\text{missed}}, \quad \lambda = 100
\end{equation}

% ---------------------------------------------------------------
\section{실험 설정}

\textbf{시나리오.} 모든 시나리오에 작업자 오류가 있다. 학습 날짜와 시험 날짜는 겹치지 않는다.
\begin{itemize}
  \item 학습: 18일 $\times$ 시드 42--43 (36개). 가벼운 날 6일(01-13, 02-10, 03-09, 05-12, 06-15, 07-31), 중간 8일(01-09, 01-12, 02-09, 04-24, 05-23, 06-26, 07-11, 10-13), 무거운 날 4일(01-25, 03-22, 05-04, 05-16). % logs/train_scenarios_fault_days18
  \item 시험: 학습에 쓰지 않은 6일 $\times$ 시드 54--56 (18개). 부하가 학습 날짜와 비슷하게 퍼지도록 골랐다(표~\ref{tab:testdays}). % logs/test_scenarios_newdays
  \item 경량 시험은 시나리오마다 5회 실행한다(90 에피소드). Webots 시험은 \todo{Webots 시험 범위: 날짜, 시드, 방법}.
\end{itemize}

\begin{table}[t]
\centering
\caption{시험 날짜. 상자 수와 주문 수는 한 시간 에피소드 기준.} % scratchpad day_loads.txt
\label{tab:testdays}
\small
\begin{tabular}{llrr}
\toprule
부하 & 날짜 & 상자 & 주문 \\
\midrule
가벼움 & 2023-01-24 & 130 & 39 \\
 & 2023-06-09 & 139 & 35 \\
중간 & 2023-06-14 & 491 & 72 \\
 & 2023-01-05 & 585 & 66 \\
무거움 & 2023-08-22 & 922 & 182 \\
 & 2023-06-19 & 2219 & 92 \\
\bottomrule
\end{tabular}
\end{table}

\textbf{비교 방법.}
\begin{itemize}
  \item 팀 기본 WES 배정
  \item 가까운 로봇, 피킹 우선 규칙 (학습 시작값)
  \item 헝가리안 최적 매칭~\cite{kuhn1955}: 매 단계 놀고 있는 로봇과 시작 가능한 작업(같은 후보 30개)을, 로봇에서 작업 출발 칸까지 거리의 합이 최소가 되도록 짝짓는다. 작업 우선순위는 쓰지 않는다. % lib/sched_hungarian.py
  \item 제안 방법: CEM 학습 배정
\end{itemize}

\textbf{Webots.}
로봇 8대, 바퀴 물리와 충돌 회피, 통행 우선순위 규칙을 포함한다.
% 학습 날짜 시험: scratchpad webots_fault.sh, webots_more.sh (2026-11-26, 27); 새 날짜 08-22: webots_nd.sh (2026-11-27)

% ---------------------------------------------------------------
\section{결과}

\subsection{경량 시뮬레이터}

\begin{table}[t]
\centering
\caption{경량 시뮬레이터 시험 결과 (작업자 오류, 처음 보는 6일, 90 에피소드 평균). 시간은 초.} % logs/train_fault_days18_cem_r*_report.json; scratchpad newdays_screen.txt
\label{tab:lite}
\small
\begin{tabular}{lrrrr}
\toprule
방법 & 주문 & 상자 & p90 & 마감 놓침 \\
\midrule
팀 기본 WES & 393 & 309 & 818 & 3.51 \\
가까운 로봇 규칙 & 401 & 306 & 833 & 3.78 \\
헝가리안 & 400 & 308 & 836 & 3.7 \\
CEM r0 & 385 & 306 & 776 & 3.60 \\
CEM r1, r2 & \multicolumn{4}{l}{\todo{학습 중}} \\
\bottomrule
\end{tabular}
\end{table}

\begin{itemize}
  \item \todo{r0--r2 평균, 팀 WES 대비 같은 에피소드끼리의 차이 $\pm$ 표준오차, 이긴 횟수}
  \item CEM r0는 팀 기본 WES보다 주문 리드타임이 8\,초(2\%) 짧았고, p90은 818에서 776\,초로 줄었다.
  \item 헝가리안은 처음 보는 날짜에서 팀 기본 WES보다 7\,초 길었다. 특히 무거운 날(08-22)에 37\,초 나빴다.
  \item \todo{날짜별 결과: 어느 부하에서 이득이 나는지}
\end{itemize}

\subsection{Webots 검증}
\label{sec:webots}

\begin{table}[t]
\centering
\caption{Webots 결과 (작업자 오류, 처음 보는 날짜). 시간은 초. 마지막 열은 같은 시나리오끼리 잰 주문 리드타임 차이의 평균 $\pm$ 표준오차와, 팀 기본 WES보다 나았던 시나리오 수.}
\label{tab:webots}
\small
\begin{tabular}{lrrrl}
\toprule
방법 & 주문 & 상자 & p90 & 팀 WES 대비 \\
\midrule
팀 기본 WES & \multicolumn{4}{l}{\todo{Webots 새 날짜 시험}} \\
헝가리안 & \multicolumn{4}{l}{\todo{}} \\
CEM & \multicolumn{4}{l}{\todo{}} \\
\bottomrule
\end{tabular}
\end{table}

\begin{itemize}
  \item \todo{Webots 결과 해석: 경량 시뮬레이터와 방향이 같은지, 헝가리안과의 차이}
  \item \todo{상자 단위 리드타임이 같은지 확인. 같다면 개선은 처리량이 아니라 주문을 끝내는 순서에서 나온다}
\end{itemize}

\subsection{학습 날짜에 대한 과적합}
\label{sec:overfit}
시드는 입고와 작업자 오류만 바꾸고 출고 주문은 바꾸지 않는다.
따라서 학습한 날짜를 시드만 바꿔 시험하면, 학습한 배정은 이미 본 주문 패턴에서 평가된다.
표~\ref{tab:overfit}은 학습 날짜 수와 시험 날짜에 따라 팀 기본 WES 대비 개선이 어떻게 달라지는지 보여 준다.

\begin{table}[t]
\centering
\caption{학습과 시험 날짜에 따른 팀 기본 WES 대비 주문 리드타임 차이 (초, 음수가 개선). 경량 값은 같은 에피소드끼리의 차이 평균으로, 학습 시드 3개의 평균이다. 단 2일 학습의 처음 보는 6일, 3일 학습의 같은 날짜, Webots의 2일 학습 값은 학습 시드 r0 하나이다.} % logs/train_fault_base_cem_r*, train_fault_base3_cem_r*, train_fault_days18_cem_r* reports; scratchpad newdays_screen.txt, agg2.py, agg_nd.py
\label{tab:overfit}
\small
\begin{tabular}{@{}llrr@{}}
\toprule
학습 & 시험 & 경량 & Webots \\
\midrule
2일 & 같은 날짜$^\dagger$ & $-39$ & $-57$ \\
2일 & 새 6일 & $+1$ & -- \\
3일$^\ddagger$ & 같은 날짜$^\dagger$ & $-36$ & -- \\
3일$^\ddagger$ & 새 6일 & $-1$ & $+4$, $+17$$^\ast$ \\
18일 & 새 6일 & \todo{} & \todo{} \\
\bottomrule
\end{tabular}

\vspace{2pt}
\footnotesize $^\dagger$ 학습한 날짜를 시드만 바꿔 시험. $^\ddagger$ 과부하(03-22 전체) 포함. $^\ast$ 08-22 하루, 시드 54--56, 학습 시드 r0, r1.
\end{table}

\begin{itemize}
  \item 학습한 날짜에서 7\% 안팎이던 개선이, 처음 보는 날짜에서는 사라졌다. Webots에서도 같았다.
  \item 처음 보는 날짜 가운데 경량 시뮬레이터에서 차이가 가장 컸던 08-22를 Webots로 돌렸지만, 학습한 배정은 팀 기본 WES와 차이가 없었다.
  \item 학습 날짜를 18일로 늘리면 \todo{결과}.
  \item 시드만 바꾼 시험은 학습한 배정의 일반화를 과대평가한다. 날짜(주문)를 나눠 시험해야 한다.
\end{itemize}

% ---------------------------------------------------------------
\section{한계}
\begin{itemize}
  \item \textbf{시뮬레이션만 사용한다.} 실제 로봇 시연은 \todo{축소 창고 시연 추가 예정}.
  \item \textbf{오류율은 합성 가정이다.} 실측 오류율 자료가 없어 한 가지 오류 수준만 썼다.
  \item \textbf{학습된 가중치는 하나로 정해지지 않는다.} 작업 종류 특징은 차이만 의미가 있어서, 서로 다른 가중치 조합이 비슷한 성능을 낸다.
  \item \textbf{과부하에서 큰 주문을 놓칠 수 있다.} 평균 비용으로 학습하면 오래된 작업의 우선순위가 낮아져, 일부 주문이 끝까지 남을 수 있다. 2일로 학습한 배정은 Webots 과부하 시나리오에서 이 실패를 보였다. \todo{18일 학습에서도 나타나는지}
  \item \textbf{Webots 시험 규모가 작다.} 시나리오마다 1회만 실행했고, 끝내지 못한 주문 하나가 평균을 수십 초 움직인다. \todo{Webots 시나리오 수 확정 후 수정}
  \item \textbf{날짜 수가 제한적이다.} 데이터셋에서 쓸 수 있는 날짜 가운데 학습 18일, 시험 6일만 썼다.
  \item \textbf{창고 하나, 로봇 8대만 실험했다.} 더 큰 창고와 로봇 수에서의 결과는 확인하지 않았다.
  \item \textbf{경량 시뮬레이터 보정 시점.} 보정은 Webots에 비켜서기 동작을 넣기 전에 했다.
\end{itemize}

% ---------------------------------------------------------------
\section{결론}
\label{sec:conclusion}
작업자가 기록 없이 재고를 옮기는 물류창고에서, 작업-로봇 짝 점수의 가중치를 경량 시뮬레이터에서 CEM으로 학습했다.
학습한 날짜를 시드만 바꿔 시험하면 개선이 크게 부풀려졌고, 처음 보는 날짜에서는 개선이 대부분 사라졌다.
\todo{18일 학습 결과(경량, Webots) 한두 문장}
학습한 배정을 평가할 때는 시드가 아니라 주문 날짜를 나눠 시험해야 한다.
\todo{향후 과제: 더 많은 날짜와 주문 특징, 확인 작업 결정의 학습, 더 큰 창고, 축소 창고 실물 시연}

\section*{Acknowledgments}
\todo{최종본에서 작성 (익명 심사본에서는 비움)}

%% Use plainnat to work nicely with natbib.
\bibliographystyle{plainnat}
\bibliography{references}

\end{document}
